\documentclass[../main.tex]{subfiles}
\begin{document}
\chapter{Introduction}
\hrule
Are these things done? 
\begin{enumerate}
    \item Your topic, in context: what does your reader need to know to understand your thesis dissertation?
    \item The relevance of your research: how does your work fit into existing studies on your topic?
    \item Your focus and scope: what specific aspect of the topic will you address?
    \item Your questions and objectives: what does your research aim to find out, and how?
    \item An overview of your structure: what does each section contribute to the overall aim?
\end{enumerate}

\hrule

Uncertainty in the most recent climate models stems from inaccurate estimates of cloud cover and ocean heat flux \cite{meehl2020}. Clouds can significantly modify both the incoming and the outgoing radiation of the Earth, while ocean currents control the oceanic heat flux that stores close to 90\% of the Earth's energy gain inside the ocean \cite{vonschuckmann2023}. 

% For clouds, one challenge is answering how and whether rain will form. \rgc{We are NOT answering which cloud will rain or even asking this question in this thesis. It is a question I ask in my proposals, and it should not be used in this context.} 


For clouds, one challenge is answering whether enough micron-sized droplets will grow fast enough to a tenth of a millimeter. After this, they will eventually rain. The answer determines the local weather. Growth of water droplets in warm clouds from $10\mu$m to $100\mu$m is poorly understood. Although scientists believe air turbulence enhances collisional droplet growth in this size range \cite{falkovich2002,bewley2013}, no formalism connects the growth rate of $\mu$m-sized droplets to the evolution of km-sized clouds. The result is a big uncertainty in the cloud dynamics of atmospheric models \cite{liu2023}.

In the oceans, the challenge is twofold. The first concerns plankton and other marine microbes that absorb CO$_2$ from the atmosphere and greatly reduce the greenhouse gas. Some of them, like cyanobacteria and diatoms, are non-motile, with small inertia effects coming from background density fluctuations. Others, like dinoflagellates, can self-propel \cite{guasto2012}. Efficient CO$_2$ absorption requires kilometer-scale plankton colonies to survive in the upper ocean \cite{dovidio2010}. To do so, the individual planktons need to effectively capture nutrients and reproduce. Both these processes demand collisions in a constraining turbulent background flow. Thus, this problem is similar to that of water droplets colliding in the atmosphere. 


Correctly estimating energy pathways of oceanic flows is the second challenge in capturing ocean dynamics. Atmospheric winds and tides from celestial bodies like the Sun, Moon, and others generate internal waves that drive turbulent mixing of energy, temperature, salinity, and nutrients throughout the ocean abyss. The bulk of the net ocean energy is stored in long-lived, large eddies. It is the eddy-wave, eddy-eddy, and wave-wave interactions that keep the ocean churning, which preserves its vast ecology and determines the heat absorbed from the atmosphere \cite{klein2019}. These interactions span from thousands of kilometers down to a few meters, and behave differently at different length scales, making it difficult to quantify them in ocean models.

In both cases, the heart of the problem lies in the large number of length scales involved. Collisions of water droplets of micron-to-millimeter scales can have a significant influence on the kilometer-scale cloud dynamics. But it is impossible to capture zillions of collisions while simulating kilometer-scale clouds. Similarly, flow interactions in scales of a few kilometers affect ocean flows at the scale of hundreds of kilometers \cite{klein2019}. No computer, now or in the foreseeable future, can resolve all the scales simultaneously. 


This thesis focuses on the individual bottlenecks. We consider idealistic scenarios, ignore some length scales, distill the problems down to their most relevant dynamics, and analyze them. Doing so, we uncover new phenomena.
% This thesis focuses on the individual bottlenecks. We consider idealistic scenarios, ignore irrelevant length-scales, distill the problems down to their most relevant dynamics, and analyze them. Doing so, we uncover new phenomena and push the knowledge of current science. Chapters \ref{fn-cstcs}, and \ref{act-cstcs} investigate collisions of inertial and self-propelled `active' particles, respectively. They show that while collisions or caustics of inertial particles are known to happen in high-strain regions, the kind of strain is the key. Particles can avoid collisions despite encountering high strain if they originate from extensive strain regions. We also show that self-propelled particles, although inertialess, cluster and collide like inertial ones due to their activity. These findings are relevant to both cloud droplets and marine microbes. However, they are restricted to mono-dispersed particle suspensions. Chapter \ref{lck-drp} extends beyond the collision of monodisperse particles to that of polydisperse particles in turbulence and captures droplet growth dynamics in cloud turbulence. 

% Chapter \ref{trunc-euler} is an intermezzo. It deals with the Galerkin-truncated Euler equation, a system that is appealing to both fluid dynamics and statistical mechanics, and shows that not resolving all the scales in the system significantly alters system dynamics. Chapter \ref{tide-eddy} then investigates tide-eddy interactions across spatio-temporal scales. It shows that tide-eddy interactions at smaller submesoscales differ qualitatively from those at larger mesoscales. The work identifies detailed pathways for energy to cascade from internal tides to dissipative scales by forming an internal wave continuum. This can lead to a significantly different impact in capturing unresolved flow energetics in the climate models. Finally, I conclude the synopsis in Chapter \ref{conc}.


% This thesis focuses on the individual bottlenecks. We consider idealistic scenarios, ignore irrelevant length-scales, distill the problems down to their most relevant dynamics, and analyze them. Doing so, we uncover new phenomena and push the knowledge of current science.

Chapter \ref{tide-eddy} investigates tide-eddy interactions across spatio-temporal scales. Astronomical tides supply one-third of the net power injected into the global ocean. They generate internal tides that drive turbulent mixing of energy, temperature, salinity, and nutrients throughout the ocean abyss. Using idealized numerical simulations with eddy and tidal forcings, we investigate how rapidly evolving internal tides exchange energy with slowly evolving eddies across length scales. We analyze flows spanning from large-scale low (O(0.1)) Rossby numbers to small-scale high (O(1)) Rossby numbers. Low Rossby number flows have coherent vortices made up of slowly evolving balanced eddies that break at higher Rossby numbers, giving way to small-scale incoherent structures. 

\begin{center}
    \textbf{Enter challenges about flow decomposition, and the energy budgets, dispersion relation.}
\end{center}

% Upon decomposing the flow into slow and fast components, we find that nonlinear energy transfer between the components increases by 100 times across low, intermediate, and high Rossby numbers. At intermediate and high Rossby numbers, this enhanced nonlinear transfer generates a broadband internal wave continuum in the frequency spectra. The frequency spectra of low Rossby number flows, with weaker slow-fast energy transfer, show peaks only at the tidal harmonics. A significant fraction of energy in the internal wave continuum at intermediate and high Rossby numbers lies away from the dispersion lines. This highlights the importance of non-dispersive fast motions in the internal wave continuum at submesoscales. Such flows are also more efficient at transferring energy across scales, and a considerable amount of energy flux responsible for cross-scale energy transfer comes from regions with low and intermediate relative vorticity and horizontal divergence. Our findings quantify the nature of submesoscale tide-eddy interactions and point out major differences with commonly studied mesoscale tide-eddy interactions. They also demonstrate how even in the absence of near-inertial wave forcings, tide-eddy interactions alone can generate an internal wave continuum at submesoscales. This work identifies detailed pathways for energy to cascade from internal tides to dissipative scales by forming an internal wave continuum. This can lead to a significantly different impact in capturing unresolved flow energetics in the climate models. 




Chapter \ref{trunc-euler} is an intermezzo. It deals with the Galerkin-truncated Euler equation, a system that is appealing to both fluid dynamics and statistical mechanics, and shows that not resolving all the scales in the system significantly alters system dynamics.

A partial differential equation (PDE) that describes geophysical flows has infinite degrees of freedom. Numerically evolving all of them is impossible. Any numerical simulation has a finite number of grid points and therefore has finite degrees of freedom. Anyone aiming to numerically evolve the Navier-Stokes equations is obliged to evolve one of its truncated versions. Scientists investigating geophysical flows are interested in numerical solutions of the Euler equations, Navier-Stokes without viscous terms. They focus on the flows at the largest scales where the Reynolds number is absurdly large. These scales are not damped by viscosity and effectively obey inviscid dynamics. For a triply periodic system, a spectral or Galerkin truncation converts the PDE of the Euler equation into a finite set of coupled ordinary differential equations in spectral space. This system is now energy-conserving, non-integrable, and has finite degrees of freedom. It reaches an equilibrium where each degree of freedom has equal energy. This equipartition of energy gives way to a $k^2$ energy spectrum in the thermalized state. 

Thermalization appears in large eddy simulations employed to spectrally resolve geophysical flows. Hyperviscous models artificially damp the smallest scales of spectral simulations. A smaller–than–required hyperviscosity leads to thermalization. Thermalization in the Galerkin-truncated system is also challenging to numerical investigations of the finite-time blow-up problem in turbulence \cite{bustamante2012}.


Although thermalization is a consequence of Galerkin truncation, studying its properties is still useful. It is a cautionary tale for geophysical physicists and a doorway to unlock new research for mathematicians. \citet{cichowlas2005} showed that thermalization starts from the smallest scales and slowly engulfs the larger ones. During this, they claimed that the large-scale flows show a Kolmogorov-like $k^{-5/3}$ energy spectrum. A vital question is whether the dynamics of the large-scale flows care about the small scales? 
  
\filler{Add timescales for thermalization background.}

% Our work, \citet{rajarshi2026a}, finds that at the onset of thermalization, before the small scales have equilibrated, the energy spectrum follows $k^{-3}$. This is a distinct, earlier stage than the thermalization dynamics at a later stage. Through dimensional analysis, we show that this $k^{-3}$ spectrum is consistent with an effective eddy viscosity acting on the larger scales. At a later time, in partially thermalized stage, the small scales have relaxes into the $k^2$ equipartition spectrum described above, and we uncover a constant energy flux across wavenumbers. A "constant flux is normally the signature of a self-similar cascade and, by dimensional analysis, would correspond to a Kolmogorov $k^{-5/3}$ spectrum in an unbounded inertial range. In some cases, a constant flux across scales defines an inertial range. Our results show that this idea should be treated with caution. 


% Different timescales related to thermalization also behaves interestingly. Upon increasing  the spatial resolution, while numerical solutions of conventional PDEs like viscous Navier-Stokes converges to the `true' solution. The Galerkin-truncated Euler system shows a different behavior. The time for the smallest scales to thermalize remains roughly constant for different system resolutions. This means that thermalization in the truncated system is not only an inevitable consequence of the numerics, but it also cannot be postponed by increasing system resolution.


In the remaining chapters, we drift from turbulent flow behaviour to explore behavior of particles in turbulence. 
Chapters \ref{fn-cstcs}, and \ref{act-cstcs} then investigate collisions of inertial and self-propelled `active' particles, respectively. 

The flow around {very small} individual colliding particles is Stokesian, i.e., highly viscous. We, however, are interested in the particle number density at scales much larger than the particle size. For example, the Kolmogorov length scale in clouds is a few tenths of millimeters, whereas water droplets, {unless experiencing significant levels of collision,} range in {tens of} $\mu$m \cite{grabowski2013}. Therefore, cloud dynamics, even in a few meters, is governed by turbulent flows where the water droplets act like point particles with finite inertia. Plankton population in a few meters of the ocean is also similar, except that their density is comparable to {that of} the background flow \cite{pedley1992}. Still, the larger scale considered here is minuscule compared to the atmospheric scale.


Chapter  \ref{fn-cstcs} deals with caustics or collisions of finitely dense inertial particles in turbulence.

\filler{details about caustics and why this is a good metric in the continuum limit of particles}


%  While caustics are known to form in strain regions, our work \citet{rajarshi2026} shows how the kind of strain affects the outcome. Furthermore, all real particles are finitely dense, but the literature focuses on particles infinitely denser than the fluid. We study caustics in two-dimensional flows for particles denser than the fluid.
% We assume one-way coupling, i.e., that particles are too small to affect the flow and are in a dilute enough suspension so as not to influence each other. In the limit of a negligible Reynolds number, a spherical particle obeys the Gatignol-Maxey-Riley equation \cite{maxey1983,gatignol1983}. We ignore finite-size and history terms in our study.


% In point-vortex flow and in 2D turbulence, only particles that pass through high-strain regions can form caustics. Caustics form after the particle has exited the high-strain region, and higher strain levels are needed for lighter particles. Survivor particles go through high strain but do not form caustics, and our explanation for this, holding for particles of any density is as follows: particles first encounter a region of extensional strain. They eject out of these regions and into compressional strain regions along low-curvature and high-velocity paths, transiting quickly through the high-strain region without forming caustics. In contrast, caustics particles approach regions of large compressive strain with a smaller velocity. Besides, their trajectories have larger curvature, enabling the particles to sample compressional strain for a longer time, and thus to form caustics. Thus, our results explain the mechanistic reason behind the observation that caustics particles stay in large-strain regions for a long time, while those that do not spend a long time in the large-strain regions can survive caustics. \cite{batge2023}.


Next in chapter \ref{act-cstcs}, we explore caustics of self-propelled `active' particles. These particles, like planktons in the ocean despite being inertialess, can significantly detach from the flow streamlines due to their self-propulsion. We investigate, whether this self–propulsion give similar behaviour to that of inertia and lead to their caustics. 

As a co-author, my contribution to this study was to contrast the caustics of active particles with inertial ones near a point vortex through singular perturbation theory and from the continuum limit. I derived asymptotics reported below and verified them with numerics. My co-authors investigated collisions of active particles in turbulence.
% \vspace{1em}
% \hrule
% \vspace{1em}


The asymptotic inertialess limit of zero Stokes number renders inertial particles as tracers. They follow fluid streamlines without any possibility of caustics in an incompressible flow. Self-propelled particles, on the other hand, can detach from the flow owing to its internal energy despite having no inertia. This can lead to their clustering \cite{durham2013}, and even to caustics \cite{chajwa2026}. Our work, \citet{chajwa2026}, investigates caustics of such self-propelled inertialess `active' particles. We show that active caustics share a close similarity with their inertial siblings.

The particles we consider are orientable, i.e., have a preferred direction of propulsion. When the background flow stretches them in that direction, they respond by moving in that direction. Therefore, when these particles are put in regions of strong flow gradients, they can deviate from flow streamlines, as if they have inertia. We show that their motility acts as an effective mass in their evolution. Their behavior in weaker flow gradients, however, is dissimilar to inertial ones.

% The easiest demonstration of this phenomenon is to put these orientable particles in a point vortex and inspect caustics as we did for inertial particles. Active particles starting out near the center of vortices fly away from the vortex center. In this region, two rings of particles can overtake each other, leading to caustics. Just like inertial particles, their dynamics is amenable to singular perturbation near the vortex origin. Far away from the vortex, in the `outer' region, their trajectories differ from inertial particles. Instead of spiraling out forever, as inertial particles do, active particles stop drifting away from the vortex center, relax to the background flow, and become tracer-like. The radial trajectories for different activity strength $\alpha$ are compared with inner and outer solutions. Just like their inertial siblings, caustics are triggered from high-strain regions close to the vortex with time for caustics diverging beyond a finite radius. Therefore, this work shows that self-propulsion can mimic inertia for active particles, which can lead to their collisions. These collisions are often vital for marine microbes to reproduce and collect nutrients.


Studies of caustics and collisions of particles so far were restricted to mono-dispersed particle suspensions. In chapter \ref{lck-drp}, we extend this to study polydisperse particles using a hybrid Euler-Lagrange framework. 

The mixing ratio of water mass to dry air mass determines the fate of a cloud in climate models. The simplest models advect this mixing ratio as tracers and trigger precipitation when its value crosses a threshold \cite{brown2023}. In reality, only droplets of certain sizes ($>100\mu$m) fall fast enough to form rain. So, droplet mass distribution is also important. Recent models use higher moment equations of the mixing ratio to capture this, but lose accuracy by assuming empirical droplet size distributions \cite{liu2023}. The super-droplet method of \citet{shima2009} is more accurate in capturing droplet growth dynamics. Despite simplified grouping of cloud droplets and collection schemes, it recovers the droplets' mass distribution function in cloud-scale models \cite{matsushima2023}. Its use is mostly confined to LES, where turbulence intensity is lower than in realistic scenarios. A lack of first-principles derivation limits the accuracy of the estimates of this model at high Reynolds numbers.

DNSs of particle-laden turbulence evolve individual droplets and accurately capture their dynamics. But they are severely bottlenecked by the pairwise collision schemes \cite{sundaram1996}, where at each time step, each droplet has to check for collisions with its neighbors. Recent findings hint that one might be able to model particle collisions without exhaustively checking for pair interactions. In a cloud of tiny droplets, only a few, larger than the mean size, accelerate the droplet size growth. Their larger size gives them a larger mass and inertia to detach significantly from the flow streamlines. Their velocity also differs from that of the neighboring tiny droplets. These larger droplets, called lucky droplets by \citet{wilkinson2016}, collide with tiny droplets on their way, and grow even bigger. The cycle continues until the droplets become big enough to start settling under gravity and eventually rain. The tiny droplets rarely collide with each other and mostly follow the flow streamlines. It is the growth rate of lucky droplets that determines whether a cloud will have enough sedimenting droplets under gravity to rain or not. 

We derive a hybrid Euler-Lagrange model for the growth of lucky droplets in turbulence. The tiny particles are evolved as a continuum with velocity and number density fields that disallow caustics. The lucky droplets are evolved Lagrangially. Their masses grow as they eat up the tiny droplets while traversing through them. Their presence makes a significant change in small droplet distribution. Our model captures droplet growth in turbulent clouds without having to check for pair collisions at each time step. This saves significant computational time.

Observational and large-scale simulation of warm clouds shows two distinct peaks in droplet size distribution \cite{chandrakar2024}. This further validates the separation of lucky and tiny droplets. We ask how the initial size of lucky droplets affects their growth.

% Keeping the net initial mass of the lucky droplets the same, we find that the smaller lucky droplets, on average, grow faster compared to its initial size. This is true both in $2$D and $3$D, although the growth rate in $3$D is much lower. However, the bigger droplets start out with an advantage. Despite their slow growth rate relative to its size, they cross a large droplet threshold earlier than the smaller lucky droplets. This hints that while the smaller lucky droplets deplete the tiny ones more, as they are more in number, individual drops may fail to grow to the size necessary for the cloud to rain. Bigger lucky drops are fewer, and will deplete the tiny droplets less. But each of these drops is more likely to reach the threshold that triggers rain. 

% In $2$D turbulence, the lucky droplets starting out inside vortical regions grow faster than the ones starting outside. They are thrown out of strong vortices, and on their way out, they collect more tiny droplets. This highlights the importance of long-lived strong vortices in droplet growth.

Finally, I conclude the thesis in Chapter \ref{conc}.
\end{document}
